%==============================================================
%  lecons/algebre/rang.tex — Théorème du rang
%==============================================================

\lecontitle{Théorème du rang}{Algèbre linéaire — Applications linéaires en dimension finie}

%=============================================================
%  § 1 — ÉNONCÉ
%=============================================================
\secnum{navyblue}{1}{Énoncé}

\begin{tcolorbox}[thmbox]
  \textbf{Théorème du rang (théorème fondamental de l'algèbre linéaire).}\enspace%
  \index{théorème du rang}\index{noyau}\index{image}
  \textit{Soit $f : E \to F$ une application linéaire, avec $E$ de dimension
  finie $n$. Alors $\mathrm{Im}\,f$ est de dimension finie et}
  \[
    \dim(\ker f) + \dim(\mathrm{Im}\,f) = \dim(E).
  \]
  \textit{En notation matricielle, pour $A \in \mathcal{M}_{m,n}(\mathbb{K})$ :}
  \[
    \dim(\ker A) + \mathrm{rg}(A) = n.
  \]
\end{tcolorbox}

\rubriq{Vocabulaire}
$\dim(\ker f)$ est la \emph{nullité}\index{nullité} de $f$ ;
$\dim(\mathrm{Im}\,f) = \mathrm{rg}(f)$ est le \emph{rang}\index{rang} de $f$.
Le théorème affirme que nullité + rang = dimension de l'espace de départ.

\decsep{}

%=============================================================
%  § 2 — DÉMONSTRATION
%=============================================================
\secnum{navyblue}{2}{Démonstration}

\rubriq{Idée}
On complète une base du noyau en une base de $E$, et on montre que les images
des vecteurs ajoutés forment une base de $\mathrm{Im}\,f$.

\begin{mdframed}[style=proofbar]
  Soit $(e_1,\ldots,e_p)$ une base de $\ker f$ (avec $p = \dim\ker f$).
  Par le théorème de la base incomplète\index{base incomplète (théorème de la)},
  on complète en une base $(e_1,\ldots,e_p,\, e_{p+1},\ldots,e_n)$ de $E$.

  \textit{Affirmation : $(f(e_{p+1}),\ldots,f(e_n))$ est une base de $\mathrm{Im}\,f$ ;
  les deux étapes ci-dessous établissent qu'elle est à la fois génératrice et libre.}

  \smallskip
  \textit{Famille génératrice.}
  Tout $y \in \mathrm{Im}\,f$ s'écrit $y = f(x)$ pour un
  $x = \sum_{i=1}^p \lambda_i e_i + \sum_{j=p+1}^n \mu_j e_j$.
  Comme $f(e_i) = 0$ pour $i \leq p$ :
  \[
    y = f(x) = \sum_{j=p+1}^n \mu_j f(e_j).
  \]

  \smallskip
  \textit{Famille libre.}
  Supposons $\sum_{j=p+1}^n \mu_j f(e_j) = 0$.
  Alors $\sum_{j=p+1}^n \mu_j e_j \in \ker f$ ; comme $(e_1,\ldots,e_p)$ est une
  base de $\ker f$, il existe \emph{certains} scalaires $\lambda_1,\ldots,\lambda_p$
  tels que
  \[
    \sum_{j=p+1}^n \mu_j e_j = \sum_{i=1}^p \lambda_i e_i
    \quad\Longrightarrow\quad
    \sum_{i=1}^p (-\lambda_i)e_i + \sum_{j=p+1}^n \mu_j e_j = 0.
  \]
  La famille $(e_1,\ldots,e_n)$ étant libre, tous les coefficients sont nuls :
  $\mu_j = 0$ pour tout $j$.

  \smallskip
  Ainsi $\mathrm{Im}\,f$ est engendré par $n-p$ vecteurs libres, donc
  $\dim(\mathrm{Im}\,f) = n - p$, et $p + (n-p) = n$.
  $\blacksquare$
\end{mdframed}

\decsep{}

%=============================================================
%  § 3 — COROLLAIRES FONDAMENTAUX
%=============================================================
\secnum{navyblue}{3}{Corollaires}

\begin{tcolorbox}[thmbox]
  Soit $f : E \to F$ linéaire, $\dim E = \dim F = n < +\infty$.
  \begin{enumerate}
    \item $f$ injective $\Leftrightarrow$ $f$ surjective $\Leftrightarrow$ $f$ bijective.
    \item $f$ injective $\Leftrightarrow$ $\ker f = \{0\}$ $\Leftrightarrow$ $\mathrm{rg}(f) = n$.
    \item \textbf{Système $Ax = b$, $A \in \mathcal{M}_n(\mathbb{K})$ :}
          le système a une unique solution $\Leftrightarrow$ $\det(A)\neq 0$.
    \item \textbf{Système $Ax = b$, $A \in \mathcal{M}_{m,n}(\mathbb{K})$ :}
          l'ensemble des solutions est soit vide, soit un sous-espace affine
          de dimension $n - \mathrm{rg}(A)$.
  \end{enumerate}
\end{tcolorbox}

\rubriq{Preuve du corollaire 1}

\begin{mdframed}[style=proofbar]
  $f$ injective $\Leftrightarrow$ $\ker f = \{0\}$ $\Leftrightarrow$ $\dim\ker f = 0$
  $\Leftrightarrow$ $\mathrm{rg}(f) = n = \dim F$ $\Leftrightarrow$ $f$ surjective. \hfill$\blacksquare$
\end{mdframed}

\rubriq{Rang et systèmes linéaires}

\begin{mdframed}[style=proofbar]
  Le système $Ax = b$ est compatible $\Leftrightarrow$ $b \in \mathrm{Im}\,A$
  $\Leftrightarrow$ $\mathrm{rg}(A|b) = \mathrm{rg}(A)$ (ce critère de compatibilité,
  parfois attribué à Rouché-Fontené\index{Rouché-Fontené (critère de)}, est hors-programme
  strict en prépa française et n'est mentionné ici qu'à titre indicatif).
  L'espace des solutions de $Ax = 0$ est $\ker A$, de dimension $n - \mathrm{rg}(A)$
  par le théorème du rang. Les solutions de $Ax = b$ forment un translaté de $\ker A$.
  \hfill$\blacksquare$
\end{mdframed}

\decsep{}

%=============================================================
%  § 4 — APPLICATIONS, CONTRE-EXEMPLES, EXERCICES
%=============================================================
\secnum{exgreen}{4}{Applications, pièges et exercices}

\rubriq{Applications}

\begin{mdframed}[style=exbar, nobreak=true]
  \textbf{1. Calcul du rang d'une matrice.}\enspace%
  Pour $A = \begin{pmatrix}1&2&3\\2&4&6\\1&1&1\end{pmatrix}$ :
  $L_2 \leftarrow L_2-2L_1$ et $L_3\leftarrow L_3-L_1$ donnent
  $\begin{pmatrix}1&2&3\\0&0&0\\0&-1&-2\end{pmatrix}$, donc $\mathrm{rg}(A)=2$
  et $\dim\ker A = 3 - 2 = 1$.

  \smallskip\noindent
  \textbf{2. Nombre de solutions d'un système.}\enspace%
  Un système de $m$ équations à $n$ inconnues avec $\mathrm{rg}(A)=r$ :
  si compatible, l'espace des solutions est de dimension $n-r$
  (une infinité de solutions si $r < n$, unique si $r = n$).

  \smallskip\noindent
  \textbf{3. Projecteurs.}\enspace%
  Si $p : E \to E$ est un projecteur ($p^2 = p$), alors $E = \ker p \oplus \mathrm{Im}\,p$
  et le théorème du rang donne $\dim E = \dim\ker p + \dim\mathrm{Im}\,p$.
\end{mdframed}

\vspace{6pt}
\rubriq{Pièges}

\begin{mdframed}[style=cexbar]
  \textbf{Le théorème porte sur $\dim E$, pas $\dim F$.}\enspace%
  Pour $f : \mathbb{R}^3 \to \mathbb{R}^5$, le rang est au plus $3$ (pas $5$),
  et $f$ ne peut pas être surjective.

  \smallskip\noindent
  \textbf{En dimension infinie, le théorème est faux.}\enspace%
  Soit $f : \ell^2 \to \ell^2$, $(x_1,x_2,\ldots) \mapsto (0,x_1,x_2,\ldots)$
  (shift). $f$ est injective ($\ker f = \{0\}$) mais pas surjective
  ($(1,0,0,\ldots)\notin\mathrm{Im}\,f$). Ici le rang n'est pas fini : l'égalité
  du théorème du rang n'a plus de sens en dimension infinie.

  \smallskip\noindent
  \textbf{$\mathrm{rg}(A) = \mathrm{rg}(A^T)$, mais $\ker A \neq \ker A^T$.}\enspace%
  Le rang-ligne et le rang-colonne coïncident, mais les noyaux de $A$ et $A^T$
  vivent dans des espaces différents ($\mathbb{K}^n$ et $\mathbb{K}^m$).
\end{mdframed}

\vspace{6pt}
\exolabel

\begin{mdframed}[style=exobar]
  \begin{enumerate}
    \item Soit $f : \mathbb{R}^4 \to \mathbb{R}^3$ définie par la matrice
          $A = \begin{pmatrix}1&2&1&0\\-1&1&0&2\\2&3&1&-2\end{pmatrix}$.
          Déterminer $\mathrm{rg}(A)$, une base de $\ker A$ et une base de $\mathrm{Im}\,A$.

    \item Soit $f : E \to E$ linéaire avec $f^2 = f$ (projecteur).
          Montrer que $E = \ker f \oplus \mathrm{Im}\,f$.

    \item Soit $A \in \mathcal{M}_n(\mathbb{K})$. Montrer que
          $\mathrm{rg}(A^2) \geq 2\,\mathrm{rg}(A) - n$
          (inégalité de Sylvester\index{Sylvester (inégalité de)}).

    \item Soit $f : E \to F$ et $g : F \to G$ linéaires. Montrer que
          $\mathrm{rg}(g\circ f) \geq \mathrm{rg}(f) + \mathrm{rg}(g) - \dim F$.

    \item Soit $A \in \mathcal{M}_{m,n}(\mathbb{R})$. Montrer que
          $\mathrm{rg}(A^T A) = \mathrm{rg}(A)$, et en déduire que
          $A^T A$ est inversible si et seulement si $A$ est injective.
  \end{enumerate}
\end{mdframed}

\leconfooter{G.\ Strang, \textit{Linear Algebra} — fondements de l'algèbre linéaire moderne}
